\documentclass[12pt]{article}
\usepackage[utf8]{inputenc}
\usepackage[T1]{fontenc}
\usepackage{hyperref}
\title{Temein Language Resources for K-12 Curriculum Development}
\author{Compiled by Claude Sonnet 5 (claude-sonnet-5) for Andr\'as Kornai}
\date{}
\begin{document}
\maketitle

\section{Introduction}
Temein (also Ron(g)e; ISO 639-3 teq) is spoken in the Nuba Hills of Sudan, with about 13,000 speakers (2006 figure) and a UNESCO rating of severely endangered. Its genealogical classification is disputed: Glottolog does not accept the Eastern Sudanic or Nilo-Saharan placement.

\section{Primary Educational Resources}
None found. There is no evidence of school materials or of instruction in Temein.

\section{Linguistic and Reference Materials}
The search found only encyclopedic and comparative sources, with no dictionary or grammar.
\begin{itemize}
\item Wikipedia article on the Temein language. \url{https://en.wikipedia.org/wiki/Temein_language}
\item Wikipedia article on the Temein languages (Temeinic: Temein and Tese). \url{https://en.wikipedia.org/wiki/Temein_languages}
\item Glottolog entry for the Temein family, with its classification note. \url{https://glottolog.org/resource/languoid/id/teme1251}
\item A Humboldt University Berlin document (edoc.hu-berlin.de) returned for the query. I did not examine it and cannot say what it contains. \url{https://edoc.hu-berlin.de/bitstreams/79bb94f0-02ac-4e19-b845-f8bf5f2c1bf8/download}
\end{itemize}

\section{Religious and Community Resources}
None found.

\section{Digital and Audio Resources}
None found.

\section{Curriculum Development Context}
Almost nothing practical is documented for Temein. A curriculum would need new descriptive work with the speech community before any materials could be written. Ongoing instability in Sudan complicates fieldwork.
\end{document}